\chapter{What makes an artefact legible}
\label{ch:legible}

\section*{What this chapter is for}

Most portfolio work fails on readability rather than on substance. A reader
gives it somewhere between five and twenty minutes, decides whether it
demonstrates judgement, and moves on. This chapter is about the properties that
survive that, and they are not the properties that feel important while you are
building.

\section{Five properties}

\subsection*{One thesis}

The artefact argues one thing. A reader should be able to state it in a
sentence after the README, and that sentence should be a claim somebody could
disagree with.

\emph{``A retrieval system''} is not a thesis. \emph{``Grading the execution
trace rather than the reply text is what makes an evaluation suite survive a
prompt rewrite''} is a thesis: it is specific, it is arguable, and everything
in the repository can then be read as evidence for it.

\judgement{The single most common defect in portfolio work is that it
demonstrates competence at assembling components and takes no position at all.
Competence at assembly is assumed at this level; it is not what is being
looked for.}

\subsection*{Readable in twenty minutes}

Not small \dash{} \emph{readable}. A large repository with an entry point, a
map, and one file that is obviously the point is readable. A small one with
five equally plausible starting places is not.

Practically: the README says what it is, what it claims, and which two files to
open. \judgement{Naming the two files is the highest-value sentence in a
portfolio README and almost nobody writes it.}

\subsection*{Gated, and the gates watched failing}

An artefact that claims correctness and has no checks is asking to be trusted,
which is what the artefact existed to avoid. An artefact with checks that have
never failed is worse, because it looks rigorous.

The discipline worth importing, and it is the estate's own: \emph{a check nobody
has watched fail is not known to be checking anything}. Introduce the defect,
watch the check refuse it, then keep the check. If you can say in an interview
that you deliberately broke your own system to confirm the suite noticed, and
that one of the breakages \emph{survived}, you have said something no
certificate says.

\subsection*{Deployed, or honest about not being}

\reported{Deployed and clickable beats a notebook.}{field-guide-get-hired} A URL
removes an entire category of doubt about whether the thing runs outside the
author's machine.

Where deployment is genuinely not worth it, say so and say what that leaves
untested. That is not an apology; it is the property in the next section.

\subsection*{It states what it does not do}

This is the one that separates a portfolio from a product demo, and it is
important enough to have its own chapter (Chapter~\ref{ch:ledger}).

\section{The README, as a structure}

\begin{kitmethod}
\textbf{One sentence: what this is.} Not the stack. The thesis.

\textbf{One paragraph: why it exists.} The problem it is a response to.

\textbf{What it does not do,} with a pointer to the register where the full list
lives. Near the top, not buried at the bottom.

\textbf{The two files to read,} named, with one line each on why those two.

\textbf{How to run it,} in commands that work in a fresh clone.

\textbf{What the checks check,} and \dash{} if you have them \dash{} what they
have caught.
\end{kitmethod}

Everything else is optional. \judgement{A badge row is worth exactly one line
and only if the badges are live; a dead badge is a claim that has stopped being
true in the most visible place on the page, which is a poor first impression in
a portfolio about rigour.}

\section{What makes an artefact read as junior}

Worth naming, because each of these is a default rather than a decision.

\begin{kittable}{@{}lX@{}}
\textbf{Signal} & \textbf{What a reader infers} \\
\midrule
A notebook as the entry point & this was explored, not built \\
No tests, or tests that assert nothing & the author has not been hurt by this yet \\
Every option configurable & no decisions were made \\
A README of screenshots & the artefact is the demo, not the reasoning \\
Generated boilerplate left in & the author did not read their own repository \\
A claim in the README the code does not support & everything else in the README is now unverified \\
\end{kittable}

The last row compounds. \judgement{One stale claim does not cost you one claim;
it costs you the reader's willingness to take the others on trust, which was
the entire asset.} Chapter~\ref{ch:claims} is the routine that prevents it.

\section{The defence is part of the artefact}

You will be asked about it, and \reported{in take-home assessments the defence
round is reported to matter more than the code
alone}{field-guide-take-home}.

So the artefact includes the answers to: why this and not the obvious
alternative; what you would do differently with a week; what in it you are least
confident about; and what it would take to run it at a hundred times the scale.
That last one is a Staff question wearing a systems costume \dash{} it is asking
whether you know which part breaks first.

\begin{drill}
Give somebody your artefact's README and twenty minutes, and ask them for three
things: what it claims, which two files they would open, and one thing they do
not believe. If they cannot answer the first two, the defect is the README. If
they have nothing for the third, the artefact is probably not claiming anything.
\end{drill}

\section*{What this chapter does not claim}

The five properties are a considered position, not a measured one. What is
reported is narrower: that deployed systems are preferred to notebooks, and
that the defence round carries weight. The rest is the author's judgement about
why some portfolio work reads as senior and most does not, and you should weigh
it as that.
